% ===================================================================
%  TABLEAU N° 1 — L'École. — Le Lycée. — La Classe.
%  Feuillets 9 a 13 du fac-simile = folios 7 a 11.
%  Correspondance : folio imprime = numero de feuillet - 2.
%
%  Les cles « %%K » sont les MEMES que dans texto/io : c'est par elles
%  que la page de lecture met les deux textes en regard. Le francais ne
%  connait ni les intitules « La bona lernanti », « La mala lernanti »,
%  « Desegno e muziko », « La doceyo pri la cienci », ni la note sur les
%  prenoms latinises : ces cles-la restent donc sans vis-a-vis, et la
%  colonne de droite y montre une case vide plutot qu'un texte deplace.
%  C'est un ecart REEL entre les deux editions, non une lacune du
%  releve : Guignon a decoupe en sections un texte que Rochelle laisse
%  courir.
%
%  ATTENTION — LE FAC-SIMILE FRANCAIS EST UNE PRISE DE VUE, non un
%  passage au scanner : l'echelle change d'une page a l'autre (de 990 a
%  1550 px de justification pour le meme bloc imprime). La calibration
%  metrique de kalibro-fr.tex est donc une MEDIANE, et le controle 11
%  devra normaliser chaque page sur sa propre largeur de bloc avant de
%  comparer quoi que ce soit.
% ===================================================================

% --- Feuillet 9 (folio 7, non imprime : page d'ouverture) ----------
\begin{VUpage}[9]{}
%  L'APPARAT DU VOLUME N'EST PAS L'OUVERTURE DU TABLEAU. Ce qui
%  precede -- le titre du livret, la serie, le filet -- annonce le
%  volume entier ; ce qui suit ouvre le tableau 1. Deux blocs, donc,
%  pour que la gravure du tableau puisse se poser entre les deux.
%%K t01-apar-0 apar
\VUsaut{2.0mm}
\VUcentre{12.6pt}{120}{LIVRET EXPLICATIF}
\VUsaut{3.2mm}
\VUcentre{8.4pt}{160}{DES}
\VUsaut{2.6mm}
\VUtitre{15.6pt}{0}{Tableaux Auxiliaires Delmas}
\VUsaut{3.4mm}
\VUornamento{9pt}
\VUsaut{5.0mm}
\VUcentre{13.2pt}{90}{PREMIÈRE SÉRIE}
\VUsaut{3.0mm}
\VUfilet{16mm}
\VUsaut{5.6mm}

%%K t01-apar-1 apar
\VUtitre{15.0pt}{60}{TABLEAU N\textsuperscript{o} 1}
\VUsaut{1.4mm}
\VUtitre{11.4pt}{0}{L'École. --- Le Lycée. --- La Classe.}
\VUsaut{2.2mm}
\VUfilet{20mm}
\VUsaut{6.4mm}

%%K t01-tit-1 sub
\VUpk{10.2pt}{40}{I. Description générale.}
\VUsaut{3.0mm}

%%K t01-01-1 p
\VUblancAlinea
1. --- Ce tableau représente une \VUgras{salle de classe}\nl
dans un \VUgras{lycée.} Dans cette salle il y a huit \VUgras{tables}\textsuperscript{(18)}\nl
et huit \VUgras{bancs}\textsuperscript{(32)}. Sur chaque banc sont assis trois\nl
élèves. Le \VUgras{professeur}\textsuperscript{(7)} est sur une \VUgras{chaise}\textsuperscript{(8)}\nl
dans sa \VUgras{chaire}\textsuperscript{(6)}.

%%K t01-02-1 p
\VUblancAlinea
2. --- A gauche de la chaire se trouve une \VUgras{armoire}\nl
\VUgras{vitrée} \textsuperscript{(15)}. A droite est le \VUgras{tableau noir}\textsuperscript{(1)} avec le\nl
\VUgras{torchon}\textsuperscript{(2)} et la \VUgras{craie}\textsuperscript{(3)}.

%%K t01-02-2 p
\VUblancAlinea
Au-dessus de la chaire sont suspendus des \VUgras{ta}\cc
\VUgras{bleaux muraux} \textsuperscript{(9, 11, 12)} et une \VUgras{carte géographi}\cc
\VUgras{que}\textsuperscript{(10)}.

%%K t01-03-1 p
\VUblancAlinea
3. --- Tous les élèves ne sont pas assis à leur\nl
\VUgras{place}\textsuperscript{(17)}. Jean\textsuperscript{(4)} est au tableau; il tient le torchon\nl
et il efface les \VUgras{figures de géométrie.}
\end{VUpage}

% --- Feuillet 10 (folio 8) -----------------------------------------
\begin{VUpage}[10]{8}
%%K t01-03-2 p
\VUblancAlinea
Pierre est auprès de la chaire; il ramasse les\nl
\VUgras{copies}\textsuperscript{(5)} et les apporte au professeur.

%%K t01-04-1 p
\VUblancAlinea
4. --- Ce professeur est le professeur de \VUgras{langues}\nl
\VUgras{vivantes.} Il enseigne aux élèves à parler, à lire et\nl
à écrire une langue étrangère \VUgras{(allemand, anglais,}\nl
\VUgras{espagnol, italien} ou \VUgras{russe)}.

%%K t01-05-1 p
\VUblancAlinea
5. --- La salle de classe est très vaste et très claire.\nl
Au fond il y a une large \VUgras{fenêtre}, deux \VUgras{vasistas} \textsuperscript{(78)}\nl
et une \VUgras{porte vitrée} pour l'aérer et l'éclairer.

%%K t01-05-2 p
\VUblancAlinea
Cette salle n'est pas froide, car il y a au milieu pour\nl
la chauffer un très bon \VUgras{poêle}\textsuperscript{(46)} avec son \VUgras{tuyau}\textsuperscript{(45)}.

%%K t01-05-3 p
\VUblancAlinea
A la cloison sont fixés des \VUgras{portemanteaux}\textsuperscript{(58)};\nl
à ces portemanteaux sont suspendus les coiffures\nl
et les manteaux des élèves.

%%K t01-06-1 p
\VUblancAlinea
6. --- L'\VUgras{horloge} \textsuperscript{(63)} marque 9 heures 5 minutes.

%%K t01-06-2 p
La \VUgras{récréation}\textsuperscript{(76)} vient de finir et la classe recom\cc
mence. La récréation a lieu dans la \VUgras{cour}\textsuperscript{(75)}. Nous\nl
voyons quelques élèves qui jouent. Un \VUgras{maître}\nl
\VUgras{répétiteur}\textsuperscript{(74)} les surveille.

%%K t01-07-1 p
\VUblancAlinea
7. --- Dans la cour nous voyons encore différentes\nl
personnes; ce sont : l'\VUgras{inspecteur} \textsuperscript{(73)}, le \VUgras{provi}\cc
\VUgras{seur} \textsuperscript{(71)} et le \VUgras{censeur} \textsuperscript{(72)}.

%%K t01-07-2 p
\VUblancAlinea
L'inspecteur inspecte les écoles, le proviseur dirige\nl
le lycée, le censeur surveille les études.

%%K t01-07-3 p
\VUblancAlinea
Le quatrième personnage est le \VUgras{concierge}\textsuperscript{(77)}. Il\nl
tient sous son bras un registre pour inscrire les\nl
absents.

%%K t01-08-1 p
\VUblancAlinea
8. --- Au fond de la cour il y a les \VUgras{appareils de}\nl
\VUgras{gymnastique}\textsuperscript{(70)} et un atelier pour les \VUgras{travaux}\nl
\VUgras{manuels}\textsuperscript{(69)}.

%%K t01-08-2 p
\VUblancAlinea
A droite du tableau nous apercevons les \VUgras{classes}\nl
de \VUgras{musique} et de \VUgras{dessin.}

%%K t01-tit-3 sub
\VUsaut{5.4mm}
\VUpk{10.2pt}{40}{II. Description détaillée.}
\VUsaut{4.0mm}

%%K t01-09-1 p
\VUblancAlinea
9. --- Les élèves du premier banc à droite de la\nl
chaire sont \VUgras{Henri}, \VUgras{Étienne} et \VUgras{Charles.}
\end{VUpage}

% --- Feuillet 11 (folio 9) -----------------------------------------
\begin{VUpage}[11]{9}
%%K t01-09-2 p
\VUblancAlinea
Henri est très attentif et très travailleur; il écoute\nl
le professeur. Sur sa table il a un \VUgras{cahier} \textsuperscript{(28)}, un\nl
\VUgras{livre}\textsuperscript{(29)}, un \VUgras{dictionnaire}\textsuperscript{(30)} et un \VUgras{plumier}\textsuperscript{(31)}.

Son livre a beaucoup de \VUgras{pages}; chaque chapitre\nl
contient plusieurs \VUgras{paragraphes} et chaque \VUgras{ligne} un\nl
grand nombre de \VUgras{mots.}

%%K t01-09-4 p
\VUblancAlinea
Son voisin Étienne\textsuperscript{(24)} est zélé. Il note sur son\nl
cahier la leçon pour la prochaine fois. Il a une jolie\nl
\VUgras{écriture.}

%%K t01-09-4b p
\VUblancAlinea
Son livre est fermé et nous ne voyons que la \VUgras{cou}\cc
\VUgras{verture} \textsuperscript{(27)}. A côté de lui se trouve son \VUgras{compas}\textsuperscript{(26)}.

%%K t01-09-5 p
\VUblancAlinea
Charles\textsuperscript{(23)} tient son livre à la main; il l'ouvre pour\nl
repasser sa leçon. Il a, comme chaque élève, un\nl
\VUgras{encrier}\textsuperscript{(25)} devant lui. Il est obéissant.

%%K t01-10-1 p
\VUblancAlinea
10. --- A la table à côté sont \VUgras{Louis, Antoine} et\nl
\VUgras{Paul.} Louis récite sa \VUgras{leçon} \textsuperscript{(22)} et répond aux \VUgras{ques}\cc
\VUgras{tions} du professeur. Antoine\textsuperscript{(21)} est très distrait; il\nl
attrape les mouches au lieu de suivre la classe. Le\nl
professeur le gronde et le blâme; parfois même il le\nl
punit. Paul\textsuperscript{(20)} est un bon camarade, gentil et com\cc
plaisant. Il est très attentif.

%%K t01-11-1 p
\VUblancAlinea
11. --- Derrière Henri se trouve \VUgras{Georges} \textsuperscript{(38)}. Ce\nl
n'est pas un mauvais garçon, mais il est \VUgras{bavard,}\nl
dissipé et remuant. Sa \VUgras{légèreté} et sa \VUgras{désobéis}\cc
\VUgras{sance} occasionnent du \VUgras{désordre} dans la classe et\nl
il mérite souvent un sévère \VUgras{avertissement.} Son\nl
\VUgras{cahier de brouillon} \textsuperscript{(35)} est sale, il y fait des \VUgras{taches}\nl
\VUgras{d'encre} avec sa plume, aussi a-t-il toujours sa \VUgras{gomme}\nl
\VUgras{à effacer} \textsuperscript{(36)}; sa \VUgras{serviette} \textsuperscript{(33)} est déchirée, car il\nl
est très négligent. Pourquoi apporte-t-il son \VUgras{porte}\cc
\VUgras{crayon}\textsuperscript{(37)} à la classe de langues vivantes ?

%%K t01-11-2 p
\VUblancAlinea
Son voisin Edmond\textsuperscript{(39)} ne l'aime pas. Il est bien un\nl
peu paresseux, mais il est silencieux et il se tient\nl
tranquille. Il ne bavarde pas ; seulement, au lieu de\nl
suivre, il préfère lire les \VUgras{histoires} de son \VUgras{livre de}\nl
\VUgras{lecture.}

%%K t01-11-3 p
\VUblancAlinea
Le troisième élève de ce banc est \VUgras{Ludovic}\textsuperscript{(40)}. Il\nl
est appliqué. Il écrit sur une feuille de \VUgras{papier}\nl
blanc. Devant lui, il a déposé son \VUgras{buvard} \textsuperscript{41)}.
\end{VUpage}

% --- Feuillet 12 (folio 10) ----------------------------------------
\begin{VUpage}[12]{10}
%%K t01-12-1 p
\VUblancAlinea
12. --- Les élèves du deuxième banc, à gauche du pro\cc
fesseur, sont très jeunes. Le premier, le petit \VUgras{Émile,}\nl
ne sait pas encore bien écrire, il apporte son \VUgras{ar}\cc
\VUgras{doise}\textsuperscript{(42)}, son \VUgras{crayon d'ardoise} et son \VUgras{éponge}\textsuperscript{(43)}.

%%K t01-12-2 p
A côté de lui sont \VUgras{Auguste} et \VUgras{Bernard}\textsuperscript{(44)}. Celui-\nl
ci travaille bien, mais sa \VUgras{tenue} est très négligée.

%%K t01-13-1 p
\VUblancAlinea
13. --- Derrière eux se trouvent trois \VUgras{paresseux.}\nl
\VUgras{Albert} n'apprend pas ses leçons. \VUgras{Jacques}\textsuperscript{(47)} dort\nl
au lieu d'écouter. Sa \VUgras{paresse} lui attire des \VUgras{répri}\cc
\VUgras{mandes} et même des \VUgras{punitions.} Enfin \VUgras{Christian}\textsuperscript{(48)}\nl
est trop léger. Il se conduit mal et ne fait aucun effort\nl
pour se corriger.

%%K t01-14-1 p
\VUblancAlinea
14. --- Leurs voisins, au contraire, sont sages.\nl
\VUgras{Ernest} \textsuperscript{(51)} aime l'ordre et la régularité, il se sert\nl
toujours de sa \VUgras{règle} \textsuperscript{(50)} et de son \VUgras{crayon} \textsuperscript{(49)}. Il est\nl
\VUgras{externe.}

%%K t01-14-2 p
\VUblancAlinea
\VUgras{François}\textsuperscript{(52)} est \VUgras{interne}; il porte l'\VUgras{uniforme} du\nl
lycée; il y mange et il y couche. C'est un bon \VUgras{cama}\cc
\VUgras{rade.}

%%K t01-14-3 p
\VUblancAlinea
Maurice taille son \VUgras{crayon} avec son \VUgras{canif}\textsuperscript{(56)}, il est\nl
très soigneux.

%%K t01-14-4 p
\VUblancAlinea
Il apporte tous ses livres, sa \VUgras{grammaire}\textsuperscript{(55)}, son\nl
\VUgras{carnet de notes}\textsuperscript{(54)}, et même un \VUgras{sous-main} \textsuperscript{(53)}.\nl
Son \VUgras{sac de classe}\textsuperscript{(57)} est par terre derrière lui.

%%K t01-14-5 p
\VUblancAlinea
Les deux tables au fond de la classe sont occupées\nl
par de \VUgras{bons élèves}\textsuperscript{(19)}.

%%K t01-15-1 p
\VUblancAlinea
15. --- Dans la \VUgras{classe de dessin}, il y a de jolis\nl
\VUgras{modèles} accrochés au mur et une \VUgras{lampe}\textsuperscript{(62)}, avec un\nl
\VUgras{abat-jour}, suspendue au plafond. Le \VUgras{professeur}\textsuperscript{(60)}\nl
est très patient, il corrige les \VUgras{dessins}\textsuperscript{(61)} des écoliers\nl
et leur montre à dessiner un \VUgras{buste} \textsuperscript{(59)} d'après\nl
nature.

%%K t01-16-1 p
\VUblancAlinea
16. --- La \VUgras{classe de musique} semble bien intéres\cc
sante. On y apprend à lire la \VUgras{musique}, à solfier les\nl
\VUgras{notes}\textsuperscript{(67)} de la gamme, écrites sur la \VUgras{portée} \textit{(do, ré,}\nl
\textit{mi, fa, sol, la, si)}, à reconnaître les \VUgras{clefs} et à chanter\nl
de jolies \VUgras{mélodies.} On pose la musique sur le
\parplein
\end{VUpage}

% --- Feuillet 13 (folio 11) ----------------------------------------
\begin{VUpage}[13]{11}
%%K t01-16-1 p suite
\VUcontinue
\VUgras{pupitre}\textsuperscript{(65)}. Un des élèves joue du \VUgras{piano} \textsuperscript{(64)} et le\nl
\VUgras{professeur de musique}\textsuperscript{(66)} bat la \VUgras{mesure}\textsuperscript{(68)}.

%%K t01-17-1 p
\VUblancAlinea
17. --- Nous apercevons un coin de l'\VUgras{atelier} de tra\cc
vaux manuels\textsuperscript{(69)}, où les enfants peuvent apprendre à\nl
raboter le bois et à limer le fer. Il n'y en a pas dans\nl
tous les lycées.

%%K t01-18-1 p
\VUblancAlinea
18. --- Le professeur de mathématiques enseigne\nl
aux élèves la \VUgras{géométrie}, l'\VUgras{arithmétique}, le\nl
\VUgras{calcul} et le \VUgras{système métrique.} Nous voyons\nl
sur le tableau les principales figures de géométrie :\nl
un \VUgras{cercle} \textsuperscript{(a)}, un \VUgras{carré}\textsuperscript{(b)}, un \VUgras{triangle}\textsuperscript{(c)} et une\nl
\VUgras{ligne} \textsuperscript{(d)}.

%%K t01-18-2 p
Nous voyons aussi une \VUgras{opération}; ce n'est ni une\nl
\VUgras{addition}, ni une \VUgras{soustraction}, ni une \VUgras{multipli}\cc
\VUgras{cation}, c'est une \VUgras{division} \textsuperscript{(e)}.

%%K t01-19-1 p
\VUblancAlinea
19. --- Sur le tableau du système métrique, nous\nl
distinguons : un \VUgras{mètre} \textsuperscript{(a)}, une \VUgras{balance} \textsuperscript{(b)} et ses\nl
\VUgras{poids}, un \VUgras{litre} \textsuperscript{(e)}, un \VUgras{décalitre}\textsuperscript{(f)}, un \VUgras{décilitre}\nl
et un \VUgras{mètre cube}\textsuperscript{(d)}.

%%K t01-19-2 p
\VUblancAlinea
Les élèves étudient aussi l'\VUgras{histoire naturelle} \textsuperscript{(11)}\nl
(la zoologie \textsuperscript{(a)}, la botanique \textsuperscript{(b)}, la géologie \textsuperscript{(c)}) et l'\VUgras{his}\cc
\VUgras{toire} \textsuperscript{(12)}.

%%K t01-19-3 p
\VUblancAlinea
Dans l'armoire se trouvent les instruments de \VUgras{phy}\cc
\VUgras{sique} \textsuperscript{(15)} et de \VUgras{chimie} \textsuperscript{(16)}.

%%K t01-19-4 p
\VUblancAlinea
Sur l'armoire sont : une \VUgras{sphère}\textsuperscript{(14)}, un \VUgras{cône} et un\nl
\VUgras{globe terrestre}\textsuperscript{(13)}.

%%K t01-19-5 p
\VUblancAlinea
Au-dessus des tableaux est suspendue une \VUgras{carte}\nl
\VUgras{géographique} représentant les cinq parties du\nl
monde : l'\VUgras{Europe}\textsuperscript{(g)}, l'\VUgras{Asie}\textsuperscript{(e)}, l'\VUgras{Afrique}\textsuperscript{(f)}, l'\VUgras{Amé}\cc
\VUgras{rique}\textsuperscript{(ab)} et l'\VUgras{Océanie}\textsuperscript{(h)}, et les deux grands Océans,\nl
l'\VUgras{Océan Pacifique}\textsuperscript{(c)} et l'\VUgras{Océan Atlantique}\textsuperscript{(d)}.

\VUsaut{6.0mm}
\VUfilet{22mm}
\end{VUpage}
